\section{Introduction}
\label{sec:intro}

An FPGA is a chip made of many small programmable logic blocks joined by programmable wires; a design is loaded onto it by setting those parts. A designer writes the design in a hardware description language (HDL), and a chain of CAD programs turns that text into a pattern on the chip. Two steps of the chain decide most of the result. Packing groups small pieces of logic into tile-sized groups, and placement decides where on the chip each group sits. Packing and placement see the design as a long list of gates and wires. The grouping the designer wrote down (this register, that adder, that block of the source) is mostly gone by then.

This paper asks one question: does keeping the designer's grouping, or any other cheap change to packing and placement, make the result better, and how can we find out honestly? The answer is a closed negative one. We answer it with hypotheses. Each hypothesis is stated in plain language, an experiment is run, and the experiment either confirms it or does not. Ideas that the experiments confirm are in the main text; ideas that failed are one line in the research log (Appendix~\ref{sec:log}). The work ends here, with a negative-results paper whose lasting contribution is a set of designs with a known best grouping, because the experiments that followed it found no change that beats noise.

\pp{The paper in four sentences.} First, we measured how much room a packer has by building designs whose best grouping is known; the standard packer leaves room in tile count and wire. Second, we tried every change we could define from information the packer can see (the designer's registers, refreshed timing weights, a starting-atom rule, a two-hop gain term, spare room in critical groups, duplication after placement); none gave a speed, wire or tile-count gain that stands apart from noise (Table~\ref{tab:tried}). Third, the changes that looked like gains were a flaw in the benchmark and two selections on the scored number. Fourth, a modern packer is near what local, answer-free information reaches, and the 2003 gains of Schabas and Brown do not replicate on VTR~9 (Section~\ref{sec:summary}).

\pp{What good means.} A real design is judged on three things at once: the chip area it needs (low), its speed (high) and its power (low). We measure area and speed. Power is not measured in any experiment of this paper. Speed means the slowest path of the design and also the paths nearly as slow, because speeding up the slowest path alone only promotes the next one (Section~\ref{sec:method}).

\pp{The data behind it.} M\'{\i}mir~III is the dataset this work feeds. Every piece of the circuit carries a record of the HDL text it came from. M\'{\i}mir stores that record beside where each piece was placed and how fast its paths are, so the data can show \emph{why} things ended up where they did, and not only where. Section~\ref{sec:mimir} says what M\'{\i}mir holds today and what it lacks.

\pp{The innovations, in plain English.}
\begin{enumerate}
\item \textbf{A packer that knows the grouping.} It is a copy of the standard open-source packer with a bonus for putting pieces of one HDL register or bus into the same tile (Section~\ref{sec:i-packer}).
\item \textbf{A yardstick with a known best grouping.} We build designs out of blocks that each exactly fill one tile, so the best tile grouping is known, and measure how far a packer is from it (Section~\ref{sec:i-yardstick}).
\item \textbf{A leak found by the yardstick.} The first test of the new packer showed a gain of nearly ten percent in wire used. The yardstick showed that the gain came from an input file that already contained the answer: a leak. With groupings that do not mirror the answer, the packer used more wire at every setting we tried (Section~\ref{sec:i-leak}). We state this plainly because it is a lesson about method: a gain that exists only when the tool is handed the solution is not a gain.
\item \textbf{A count of what the standard packer keeps together on slow paths} (Section~\ref{sec:i-census}).
\item \textbf{A negative baseline.} Starting placement from the designer's grouping instead of the tool's own start did not help (Section~\ref{sec:i-seed}).
\item \textbf{A second yardstick from real designs.} The synthetic yardstick is too regular: its wiring hardly changes with the size of the region around a block, as real wiring does (Section~\ref{sec:i-rent}). So we edit real designs until their grouping is the known best (Section~\ref{sec:i-yardstick-real}).
\item \textbf{Where the greedy packer goes wrong.} We measure the connections the packer treats as nearly the slowest, show that most of its picks from them are for connections that are not slow in the finished design, and keep a catalogue of the HDL patterns for which the pick goes wrong (Sections~\ref{sec:e-heap} and~\ref{sec:e-pick}).
\end{enumerate}

\pp{Contributions.} Each is an innovation with the hypotheses that the experiments settled; the headline is the CAD one (mean percent win and wins, losses and ties, in the order speed, wire, area), then the register verdict. (1) The standard packer scatters the source's registers across tiles (H8), and the new packer keeps them together without using more tiles (H9a; mechanism checks, not end results), all \textsc{useful}. (2) The yardstick is exact (H24, H27, H29--H31: \textsc{useful}) and shows room: the standard packer uses 22\% more wire than the known grouping over all topologies and 49\% more on meshes and trees (H25/H19, H32: \textsc{useful}); on 33 real designs it needs 1.108 times the proven lower bound on tiles and 1.265 times the wire of the known grouping, with the delay \textsc{inconclusive} (H38: \textsc{useful}). (3) The new packer's gain exists only when the HDL registers equal the blocks (H33, H36, H35: \textsc{useful}, an upper bound and a harm finding); with realistic registers it uses 65\% to 150\% more wire and 33\% to 41\% more delay (H22$'$: \textsc{dismissed}). (4) Where the standard packer goes wrong is measured: 60\% of its tile-count gap is a pick of an atom of another known group while the right atom was queued (Q-N2), and no feature computable from the circuit alone ranks the right atom well once the construction is removed (Q-N3). (5) The deployable changes all failed to beat noise. Speed: refreshing the timing weights $-0.5\%$ (H-refresh, \textsc{inconclusive}), a two-hop gain term $+0.7\%$ with 17 wins, 14 losses, 4 ties (\textsc{inconclusive}), spare room $-1.4\%$ to $-2.1\%$ and duplication $-0.5\%$ and $-0.4\%$ (\textsc{dismissed} or \textsc{inconclusive}); each costs or gains about one percent of wire and, for spare room, 11\% to 17\% more tiles (Table~\ref{tab:tried}). (6) A starting-atom oracle that uses the answer key improves wire by 2.8\% (25 wins, 8 losses) and the tile count by 1.2\% (23, 6, 4 ties) but not speed ($-0.7\%$; 13 wins, 20 losses), and rebuilds only 25.5\% of the known groups (Q-N1b). (7) A sweep of tool options gives eight \textsc{possibly good} area and speed effects (Section~\ref{sec:e10}), and the packer's idea of which connections are slowest is stale for 81\% of its picks (HP-1, a measurement). In total the research log has 88 rows: 15 \textsc{useful}, 10 \textsc{possibly good}, 27 \textsc{inconclusive}, 25 \textsc{dismissed}, and eleven measurements; ten untested ideas, listed in Section~\ref{sec:future}, sit beside them.
